% Generated draft: requires mathematical and source review.
\documentclass[11pt]{article}
\usepackage[margin=25mm]{geometry}
\usepackage{fontspec}
\setmainfont{Latin Modern Roman}
\usepackage{amsmath,amssymb,mathtools,unicode-math}
\setmathfont{Latin Modern Math}
\usepackage{enumitem,needspace,xurl,hyperref}
\setlength{\parindent}{0pt}
\setlength{\parskip}{5pt}
\setlength{\emergencystretch}{3em}
\newcommand{\N}{\mathbb{N}}
\newcommand{\Z}{\mathbb{Z}}
\newcommand{\Q}{\mathbb{Q}}
\newcommand{\R}{\mathbb{R}}
\newcommand{\C}{\mathbb{C}}
\newcommand{\F}{\mathbb{F}}
\newcommand{\A}{\mathbb{A}}
\newcommand{\PP}{\mathbb{P}}
\newcommand{\E}{\mathbb{E}}
\newcommand{\eps}{\varepsilon}
\newcommand{\dd}{\,\mathrm{d}}
\begin{document}
\section*{UnsolvedMath ID: 30004155}
\textbf{Source problem code:} catalogue entry
\subsection{Definitions and mathematical statement}
The conventional formulation of Vizing's conjecture concerns finite simple undirected graphs. The supplied record says ``arbitrary graphs'' without explicitly specifying these conventions; the formulation below uses this conventional graph class, subject to confirmation against the source report.

For a finite simple undirected graph \(G=(V(G),E(G))\), a dominating set is a subset \(D\subseteq V(G)\) such that every vertex in \(V(G)\setminus D\) is adjacent to a vertex in \(D\). Its domination number is
\[
\gamma(G)=\min\{\lvert D\rvert:D\subseteq V(G)\text{ is a dominating set of }G\}.
\]
This definition gives \(\gamma(G)=0\) when \(V(G)=\varnothing\). For graphs \(G\) and \(H\), their Cartesian product \(G\square H\) has vertex set \(V(G)\times V(H)\). Distinct vertices \((g,h)\) and \((g',h')\) are adjacent precisely when either \(g=g'\) and \(\{h,h'\}\in E(H)\), or \(h=h'\) and \(\{g,g'\}\in E(G)\).

The conjectural assertion is
\[
\forall\,G,H\text{ finite simple undirected graphs},\qquad
\gamma(G\square H)\geq\gamma(G)\gamma(H).
\]
There are no connectivity or positive-degree assumptions in this formulation. The original supplied wording asks whether the conjecture holds, while the cleaned statement asks for a proof of this inequality for all such pairs.

The record additionally asks for symmetry-based or symmetry-reduced, or Gröbner-basis-free, sum-of-squares and semidefinite certificates for the conjecture's algebraic formulation at larger parameter values. A sum-of-squares polynomial over \(\R\) is a polynomial of the form \(\sum_{i=1}^{m}q_i^2\), where the \(q_i\) are real polynomials. However, the record does not provide the algebraic formulation, its variables, coefficient field, polynomial constraints or ideal, target polynomial, parameter definitions, or the parameter range already treated. Consequently, this additional certificate-development request cannot be made into a fully specified mathematical assertion from the supplied evidence. No particular encoding or certificate identity is assumed here.
\subsection{Short English statement}
Vizing's conjecture asserts that the domination number of the Cartesian product of two finite simple undirected graphs is at least the product of their domination numbers. The record also requests larger-parameter algebraic certificates using symmetry or methods avoiding Gröbner bases, but does not specify the algebraic encoding or parameter range. The graph conventions and the certificate target require source review.
\subsection{Sources}
\begin{itemize}
\item Supplied local UnsolvedMath registry record, permanent ID \(30004155\), title ``Vizing's Conjecture and Algebraic Sum-of-Squares Certificates,'' source problem code \texttt{catalogue entry}, in the Oberwolfach Reports Open Problems collection. The supplied proposed year is 2019; no individual proposer is identified.
\item Supplied primary report citation: \emph{Mixed-integer Nonlinear Optimization: a hatchery for modern mathematics} (2020), Oberwolfach Reports, pp.\ 1573--1637, DOI \texttt{10.4171/owr/2019/26}. \url{https://doi.org/10.4171/owr/2019/26}. This is also the supplied literature-source link labelled ``OWR primary report.'' The background describes Oberwolfach Reports as documenting mathematical workshops and their open-problem discussions.
\item Supplied literature-source link labelled ``constant-factor progress on Vizing's conjecture'': \url{https://arxiv.org/abs/2606.14414}. No authors, paper title, or precise constant-factor theorem are supplied.
\item Literature assessment checked at \texttt{2026-08-22}: the supplied curation labels the problem open and states that Vizing's conjecture remains unresolved for arbitrary graphs, that recent work proves only improved constant-factor lower bounds, and that the sum-of-squares formulation does not constitute a proof of the conjecture. This is supplied dated curation derived from the Oberwolfach Reports record, not an independently verified literature assessment; the time-sensitive status and cited claims require independent verification.
\item The registry labels the statement \texttt{corrected\_verified} and says that its mathematical content was preserved while notation, typography, and sentence structure were repaired for publication. These are supplied editorial claims, not independent mathematical verification of this draft.
\end{itemize}
\end{document}
